\scanpage{39}
\deleted{Rappelons la \unsure{caractérisation|construction}}

Notations $(\lambda_i(u))$ si $u\in L_0(E)$ : c'est la suite des valeurs \deleted{\unsure{caractéristiques}} propres \inserted{non nulles} de $u$, rangées par ordre de modules \unsure{non} \unsure{décroissants|croissants}, complétée par des $0$ si \unsure{cette} suite est finie.

Si la dim de $E$ est finie, \inserted{soit $n$} il y a lieu \inserted{d'attacher} \margin{?} son intérêt \deleted{\unsure{seulement}} à $\lambda_1\unsure{(u)},\ \unsure{\ldots}\ \lambda_n(u)$ \inserted{Si le rang est fini, \unsure{on s'arrête|on garde} même \deleted{\unsure{au terme|au rang}} à partir du rang $\unsure{n|r}+1$} On pose de plus \deleted{$\unsure{\rho|f}_n(u)\unsure{=|:}$} , pour $u\in L_0(E,F)$
\[
\boxed{\boxed{\rho_n(u)=\lambda_n(\lvert u\rvert)}}
\]
donc $\rho_n(u)$ est une \unsure{suite} \unsure{décroiss\textsuperscript{te}|décroissante|décroist} de nbs positifs.

\[
\deleted{\unsure{\rho_1(u)\leq\lVert u\rVert}}\qquad\deleted{\unsure{\lambda_1(u)}}
\]
\[
\boxed{\lvert\lambda_1(u)\rvert\leq\rho_1(u)=\lVert u\rVert}
\qquad\boxed{\rho_n(u)=\rho_n(u^*)}
\]

Si $h$ est \deleted{\unsure{hermitien}} positif, on sait que
\[
\boxed{\rho_n(h)=\mathop{\rm Inf}\limits_{E_{n-1}}
\deleted{\mathop{\rm Sup}\limits_{\unsure{x\in E}}\lVert\unsure{hx}\rVert}
\unsure{(}\mathop{\rm Sup}\limits_{\substack{x\in E_{n-1}\\\lVert x\rVert\leq1}}(hx,x)\bigr)\deleted{{}^{\unsure{1/2|2}}}}
\]
% slanted
\margin{$\unsure{\mathscr H_0|\mathscr H_\infty}$} où \deleted{$h_F$ est la restriction à $F\times F$ de la forme} \inserted{$E_{n-1}$ parcourt les} sous-espaces vectoriels fermés de codimension $\unsure{n-1|n+1|n}$. \deleted{sesquilinéaire $(hx,y)$}

On en conclut $\rho_n(h^{\unsure{4|2|a}})^2=\rho_n(h^{\unsure{2|7}})=\mathop{\rm Inf}\limits_{F_n}\mathop{\rm Sup}\limits_{\substack{x\in F_n\\\lVert x\rVert\leq1}}\lVert hx\rVert^{\unsure{2|4}}$, \deleted{\unsure{d'où}} d'où
\[
\begin{aligned}
&\deleted{\rho_n(h)=\mathop{\rm Inf}\mathop{\rm Sup}\limits_{x\in F_n}hx}\\
&\deleted{\text{d'où }\rho_n(u)^2=\rho_n(u^*u)=\mathop{\rm Inf}\limits_{\unsure{E_n}}\mathop{\rm Sup}\limits_{x\in E_n}\lVert u^*ux\rVert}\\
&\rho_n(h)\unsure{'}=\mathop{\rm Inf}\limits_{E_n}\mathop{\rm Sup}\limits_{\substack{x\in E_{n-1}\\\lVert x\rVert\leq1}}\lVert hx\rVert\deleted{{}^{\unsure{2}}}\qquad \deleted{\unsure{(h\geq0)}}\ (h\geq0)
\end{aligned}
\]
D'où, en posant $h=\lvert u\rvert$
\[
\boxed{\rho_n(u)=\mathop{\rm Inf}\limits_{F_n}\mathop{\rm Sup}\limits_{\substack{x\in E_{n-1}\\\lVert x\rVert\leq1}}\lVert ux\rVert}
\]

On \deleted{\unsure{en conclut}} \inserted{en \unsure{conclut}} aussitôt :
\[
\boxed{\begin{aligned}
\rho_n(uv)&\leq\lVert u\rVert\,\rho_n(v)\\
\rho_n(uv)&\leq\lVert v\rVert\,\rho_n(u)
\end{aligned}}
\]
Et de \unsure{m\textsuperscript{me}|même} : \uline{$\rho_n(h)$ est fonction croissante de l'hermitien positif $h$.} \deleted{\uline{Corollaire} Soit $u\in L_0(E,F)$} \unsure{Enfin|Écrivons} \deleted{\unsure{on}}
\[
\deleted{\boxed{\unsure{\sum_i}\rho_i(u)=\mathop{\rm Sup}\limits_{\unsure{p_n}}\unsure{\mathrm{Tr}}\,\lvert\unsure{p_nu}\rvert}}
\]

\margin{?} \deleted{On \unsure{en} \unsure{démontre} \unsure{Tr} \illegible{1}} les inégalités de Ky-Fan
\[
\boxed{\begin{aligned}
\rho_{m\unsure{+|-}n\unsure{-|+}1}(A\unsure{B|\Pi})&\unsure{\leq|=}\rho_{m+1}(A)\rho_{n+1}(B)\\
\rho_{m\unsure{+|-}n\unsure{+|-}1}(A+\unsure{B|\Pi})&\unsure{<|\leq}\rho_{m+1}(A)+\rho_{n+1}(B)
\end{aligned}}
\]
% archive sheet number: 40

